%---------------------------------------------
\section{Stilistische Anforderungen}
%---------------------------------------------
Die "`Ich"'- oder "`Wir"'-Perspektive soll im gesamten Textteil vermieden werden.
Stattdessen sollte aus der dritten, neutralen Person oder in Passivkonstruktionen geschrieben werden \cite[S.~130]{Hering2009}.
Die wissenschaftliche Arbeit darf nicht im Sinne eines Erlebnisberichtes sondern als sachorientierte Dokumentation aufgefasst werden. 
Eine Ausnahme besteht im Schlussteil, in dem der Autor eine kritische Auseinandersetzung mit der Thematik einbringen darf.

Ein wesentliches Kriterium für die Verständlichkeit der Inhalte ist ein wissenschaftlich-technischer, konkreter, klarer und präziser Sprachstil.
Es sollte auf keinen Umständen Umgangssprache verwendet werden. 
Folgende Punkte sorgen für bessere Textverständlichkeit:
\begin{itemize}
	\item Kurze Sätze bilden.
	\item Möglichst nur einen Hauptsatz mit einem oder zwei Nebensätzen bzw. zwei mit Semikolon verknüpfte Hauptsätze verwenden. 
	\item Nicht zu viele Nebensätze in Klammern sondern als alleinstehende Sätze oder in Fußnoten unterbringen.
	\item "`Welcher, welche, welches"' als Relativpronomen sind veraltet und sollen ersetzt werden durch "`der, die, das"'. 
	\item Ein- und überleitende Sätze verwenden, um Ihre Leser mit Worten zu führen.
Damit können Sie:\\
\begin{itemize}
	\item die Gliederung bzw. das Inhaltsverzeichnis aufgreifen,
	\item den bisher beschriebenen Sachverhalt zusammenfassen,
	\item zum nächsten Dokumentteil überleiten oder
	\item einen neuen Dokumentteil einleiten.
\end{itemize}
	\item Anschaulich formulieren, in Bildern sprechen!
	Analogien, Metaphern und Vergleiche erzeugen Assoziationen beim Leser.
	Dadurch kann der Leser Ähnlichkeiten und Abweichungen von dem ihm bereits bekannten Wissen leichter erkennen \cite[S.~129]{Hering2009}.
	\item Konjunktive vermeiden.
	\item Apostrophe sparsam im Text verwenden.
	\item Indirekte statt direkte Zitate verwenden.
	\item Wiederholungen von Wörtern in Sätzen vermeiden.
	\item Füllwörter, d.\,h. Wörter die beim Weglassen keinen Einfluss auf den Inhalt haben, vermeiden.
	\item Inhaltliche Blöcke im Text bilden und diese durch Einrückungen voneinander trennen.
\end{itemize}

\noindent
Die Erzählzeit ist das Präsens (Gegenwart).
Das Imperfekt (Vergangenheit) wird nur bei Beschreibung eines früher verwendetes Bauteil, Messverfahren usw., eingesetzt.
Es sollten keine emotionsbeladenen und umgangssprachlichen Formulierungen verwendet werden \cite[S.~131]{Hering2009}.
Zur verständliche Formulierung gehören einleitende Sätze am Anfang eines Dokumentteils.
Zusammenfassende und überleitende Sätze sollten am Ende eines Dokumentteils eingebracht werden.
Diese Sätze können auch Einsatz zur Ein- und Überleitung bei Bildern und Tabellen finden.
Sie beschreiben z.B. welche Informationen und warum als nächstes folgen, wie die folgenden Informationen zu bewerten sind, welchen Stellenwert die zuvor beschriebenen Informationen haben usw. (vgl. \cite{Hering2009}, S. 136 f.).

Technische Geräte sollen grundsätzlich nach ihrer Funktion benannt werden.
Zum Beispiel ist die Bezeichnung Antriebsrad besser als Zahnrad oder Halteplatte eindeutiger als Platte.
Technische Verfahren werden zur Veranschaulichung funktionsgerecht benannt.
So sollte die Bezeichnung "`Variante 1"' in z.\,B. "`elektrisch-mechanische Lösung"' umbenannt werden \cite[S.~131]{Hering2009}.
Dabei muss aber beachtet werden, dass keine Worte erfunden werden dürfen.
Es sollen nur Worte verwendet werden, die im Duden stehen.
Bezeichnungen wie z.B. "`bestromen"' oder "`anflanschen"' sollten in einer wissenschaftlichen Arbeit durch gebräuchliche Ausdrücke ersetzt werden.

Bei der Benutzung von Fachbegriffen, Fremdwörten und Abkürzungen gilt es eine angemessene Verwendung einzuhalten.
Fremdsprachige Wörter sollten vermieden werden.
Nach Möglichkeit sollen die sinngemäßen Bezeichnungen in der deutschen Sprache verwendet werden.
Ausnahme bilden fremdsprachige Fachausdrücke, bei denen eine Übersetzung ins Deutsche übertrieben erscheinen würde (sogenannte Anglizismen vermeiden).
Bei der Verwendung von Fachbegriffen müssen diese bei erstmaliger Verwendung erläutert werden.
Dies kann sowohl innerhalb des Textes als auch in Fußnoten erfolgen.\footnote{In den Fußnoten können Anmerkungen untergebracht werden ohne den Lesefluss zu stören.
In die Fußnoten gehören ausschließlich zusätzliche Informationen und Anmerkungen, die als nennenswert erachtet werden, aber nicht für die Argumentation erforderlich sind. Jede Fußnote beginnt mit einem Großbuchstaben und endet mit einem Punkt \cite[S.~197~f.]{Balz2008}.}
Unbekannte Abkürzungen werden ebenfalls beim ersten Auftreten definiert \cite[S.~136]{Balz2008}.
Sobald eine Abkürzung einmalig eingeführt wurde, muss sie konsequent verwendet werden.
Beim Einführen einer Abkürzung wird zuerst der vollständige Begriff genannt und in Klammern auf die Abkürzung verwiesen wird.
Bei mehr als fünf verwendeten Abkürzungen sollte ein Abkürzungsverzeichnis erstellt werden.

